\section{Entornos gráficos}

Antes de instalar un escritorio hay que asegurarse de que los drivers de vídeo
están correctamente instalados. Es uno de los puntos donde más alumnos se atascan.

\subsection{Identificar la tarjeta gráfica}

\begin{tcolorbox}[title=Identificamos el hardware gráfico]
\begin{lstlisting}[style=bashstyle]
lspci | grep -i vga
\end{lstlisting}
\end{tcolorbox}

\subsection{Drivers de vídeo}

\begin{tcolorbox}[title=Intel o AMD — Mesa (recomendado)]
\begin{lstlisting}[style=bashstyle]
yay -S mesa
\end{lstlisting}
\tcblower
\scriptsize Proporciona soporte OpenGL y Vulkan para Intel y AMD. Opción por defecto y recomendada para Wayland.
\end{tcolorbox}

\begin{tcolorbox}[title=NVIDIA — driver propietario]
\begin{lstlisting}[style=bashstyle]
yay -S nvidia nvidia-utils
\end{lstlisting}
\tcblower
\scriptsize Para RTX 2000 o más nuevas puede usarse también \texttt{nvidia-open}. Para kernels LTS, usa \texttt{nvidia-lts}.
\end{tcolorbox}

\begin{tcolorbox}[title=Hardware antiguo o genérico]
\begin{lstlisting}[style=bashstyle]
yay -S xf86-video-vesa
\end{lstlisting}
\end{tcolorbox}

\subsection{Entornos de escritorio}

\begin{tcolorbox}[title=Instalar \textbf{GNOME}]
\begin{lstlisting}[style=bashstyle]
yay -S gnome gdm
systemctl enable gdm
\end{lstlisting}
\tcblower
\scriptsize Entorno completo con excelente soporte Wayland. \texttt{gdm} es su display manager nativo.
\end{tcolorbox}

\begin{tcolorbox}[title=Instalar \textbf{KDE Plasma}]
\begin{lstlisting}[style=bashstyle]
yay -S plasma sddm
systemctl enable sddm
\end{lstlisting}
\tcblower
\scriptsize Altamente configurable. \texttt{sddm} es el display manager recomendado por KDE.
\end{tcolorbox}

\begin{tcolorbox}[title=Instalar \textbf{XFCE} (opción ligera)]
\begin{lstlisting}[style=bashstyle]
yay -S xfce4 xfce4-goodies lightdm lightdm-gtk-greeter
systemctl enable lightdm
\end{lstlisting}
\tcblower
\scriptsize Ideal para hardware con pocos recursos. Consume sensiblemente menos RAM que GNOME o KDE.
\end{tcolorbox}

\subsection{Compositores Wayland (caso avanzado)}

\begin{tcolorbox}[title=\textbf{Hyprland} — tiling compositor Wayland]
\begin{lstlisting}[style=bashstyle]
yay -S hyprland waybar wofi kitty
\end{lstlisting}
\tcblower
\scriptsize Se inicia desde la TTY con \texttt{Hyprland}. Sin necesidad de display manager.
\end{tcolorbox}

\begin{tcolorbox}[title=\textbf{Sway} — equivalente Wayland de i3]
\begin{lstlisting}[style=bashstyle]
yay -S sway swaybar dmenu foot
\end{lstlisting}
\tcblower
\scriptsize Se inicia con \texttt{sway}. Más estable y conservador que Hyprland.
\end{tcolorbox}

\subsection{Reinicio y verificación}

\begin{tcolorbox}[title=Salimos del chroot y reiniciamos]
\begin{lstlisting}[style=bashstyle]
exit
umount -R /mnt
reboot
\end{lstlisting}
\tcblower
\scriptsize Al arrancar deberías ver la pantalla de login del display manager.
Si el sistema arranca en TTY, comprueba: \texttt{systemctl status gdm} (o \texttt{sddm}, \texttt{lightdm}).
\end{tcolorbox}
